% ===========================================================================
\chapter{Drawing the map}\label{ch:drawing}

Now for a picture. Nex's \code{graphics/Canvas} library gives us a canvas
of pixels to draw on, which we can save as a PNG file or show in a window.
This chapter draws the island's land, water, woods and roads; the next one
adds the words.

\section{From metres to pixels}

The canvas does the arithmetic of turning map coordinates into pixels for
us. We tell it which rectangle of the map to show, with
\code{set_x_scale(x0, x1)} and \code{set_y_scale(y0, y1)}, and then draw
in map coordinates: metres. A canvas \code{w} pixels wide showing $x$ from
$x_0$ to $x_1$ puts the point $x$ at pixel column
\[
  \frac{x - x_0}{x_1 - x_0}\, w,
\]
and similarly for $y$, except that pixel rows count \emph{down} from the
top of the image, while our $y$ counts \emph{up} towards the north. The
canvas turns it over for us, so north is at the top, as on any map.

The one thing we must get right is the shape. If the rectangle we show
has a different shape from the canvas, the map is stretched. So instead of
choosing a rectangle, we choose a \emph{centre} and a \emph{scale}, in
metres per pixel, and work out the rectangle from those:
\[
  x_0 = c_x - \tfrac{1}{2}\, w \cdot \text{scale}, \qquad
  x_1 = c_x + \tfrac{1}{2}\, w \cdot \text{scale},
\]
and the same for $y$ with the canvas's height. Every pixel is then the same
number of metres wide as it is high. Zooming in is making the scale
smaller; panning is moving the centre.

There is one complication, for later. In Chapter~\ref{ch:livemap} the live
map draws panels across the top and the bottom of the window, with
instructions and settings, and the map should be centred in the space
between them. So the view knows how many pixels at the top and bottom are
covered, its \emph{insets}, and centres the map in the rest. For now both
are zero.

\section{The colours}

Map makers spend a long time on colours. Ours are close to the ones
OpenStreetMap's standard style uses, which are tuned so that the important
things (roads and names) stand out from the background (land and water):
pale, quiet colours for areas, stronger ones for roads.

\begin{newfile}{map_view.nex}
-- map_view.nex: drawing the map on a canvas, at any zoom.

intern Island_Map
intern Helpers
intern graphics/Canvas

-- The map's colours.
class Palette
feature
  SEA = create Color.make(170, 211, 223)
  LAND = create Color.make(242, 239, 233)
  SAND = create Color.make(240, 228, 190)
  TOWN = create Color.make(232, 224, 218)
  WOOD = create Color.make(181, 214, 160)
  WATER = create Color.make(170, 211, 223)
  RIVER = create Color.make(120, 175, 205)
  CASING = create Color.make(150, 140, 130)
  MAIN_ROAD = create Color.make(250, 205, 140)
  MAIN_CASING = create Color.make(205, 140, 70)
  MINOR_ROAD = create Color.make(255, 255, 255)
  TRACK = create Color.make(160, 115, 70)
  FOOTPATH = create Color.make(225, 95, 80)
  LABEL = create Color.make(60, 60, 60)
  ROAD_LABEL = create Color.make(90, 85, 80)
  WATER_LABEL = create Color.make(50, 110, 150)
  HALO = create Color.make_rgba(255, 255, 255, 220)
end
\end{newfile}

\code{HALO} is white and slightly see-through: its fourth number, the
\emph{alpha}, is 220 out of 255.

\section{The view}

A \code{Map_View} holds the map, a canvas, and where the canvas is looking.

\begin{append}{map_view.nex}
class Map_View
create
  -- A view of `m` on a new canvas `width` by `height` pixels, showing the
  -- whole map.
  make(m: Island_Map, width, height: Integer) do
    island := m
    canvas := create Canvas.make(width, height)
    centre := create Point.make(0.0, 0.0)
    metres_per_pixel := 1.0
    show_all()
  end
feature
  island: Island_Map
  canvas: Canvas
  -- The point of the map at the middle of the canvas...
  centre: Point
  -- ...and how many metres of the map one pixel covers. Bigger means
  -- zoomed out.
  metres_per_pixel: Real

  -- Pixels at the top and bottom of the canvas that something else covers
  -- (a program's own panels), to keep the scale bar and north arrow clear
  -- of.
  inset_top: Real
  inset_bottom: Real

  set_insets(top, bottom: Real)
  require
    non_negative: top >= 0.0 and bottom >= 0.0
    room_left: top + bottom < canvas.height.to_real
  do
    inset_top := top
    inset_bottom := bottom
    look_at(centre, metres_per_pixel)
  end
\end{append}

To show the whole map, take the map's width and height in metres, and see
how many metres per pixel it takes to fit each of them on the canvas. The
larger of the two fits both; making it 8\% larger leaves a margin of sea
all round.

\begin{continued}{map_view.nex}
  -- Zoom and pan so that the whole map fits, with a little sea around it.
  show_all() do
    let w: Real := island.max_x() - island.min_x()
    let h: Real := island.max_y() - island.min_y()
    let mpp: Real := (w / canvas.width.to_real).max(h / free_height()) * 1.08
    look_at(create Point.make((island.min_x() + island.max_x()) / 2.0,
      (island.min_y() + island.max_y()) / 2.0), mpp)
  end

  -- Centre the view on c, at mpp metres per pixel. "Centre" means the
  -- middle of the part of the canvas that the insets leave free.
  look_at(c: Point, mpp: Real)
  require
    positive_scale: mpp > 0.0
  do
    centre := c
    metres_per_pixel := mpp
    let half_w: Real := canvas.width.to_real * mpp / 2.0
    let top: Real := c.y + (inset_top + free_height() / 2.0) * mpp
    canvas.set_x_scale(c.x - half_w, c.x + half_w)
    canvas.set_y_scale(top - canvas.height.to_real * mpp, top)
  end

  -- The height of the canvas that the insets leave free, in pixels.
  free_height(): Real do
    result := canvas.height.to_real - inset_top - inset_bottom
  end

  -- A distance on the canvas in pixels, as metres on the map.
  metres(pixels: Real): Real do
    result := pixels * metres_per_pixel
  end
\end{continued}

In \code{look_at}, the centre \code{c} should land in the middle of the
free part of the canvas: \code{inset_top + free_height() / 2.0} pixels
down from the top. So the top edge of the canvas is that many pixels'
worth of metres north of \code{c}, and the bottom edge is the canvas's
whole height further south.

\code{metres} converts a size in pixels into metres at the current scale.
We use it for everything that should look the same size at every zoom:
dots, arrows, the gap between a label and its place. Roads, too, are drawn
a fixed number of pixels wide, however far we zoom, as on every map app.

\section{Layers}

A map is drawn in \emph{layers}, from the bottom up: the sea first, then
the land on top of it, then lakes and woods on the land, then rivers, then
roads, then the dots for places. Each layer paints over the ones below.

\begin{continued}{map_view.nex}
  -- Draw the whole map, layer by layer, from the sea up.
  draw_map() do
    canvas.clear(Palette.SEA)
    draw_areas()
    draw_rivers()
    draw_roads()
    draw_places()
    draw_scale_bar()
    draw_north_arrow()
  end

  draw_areas() do
    across island.areas as a do
      canvas.set_pen_color(area_color(a.kind))
      fill_polygon(island.points_of(a.node_ids))
    end
  end

  draw_rivers() do
    across island.rivers as r do
      line(island.points_of(r.node_ids), Palette.RIVER, 2.5)
    end
  end
\end{continued}

The order of the areas in the map file is the order they are drawn: the
island first, then the town, the woods and the lake on top of it. That is
why \code{area_at} in the last chapter took the \emph{last} area containing
a point.

\section{Roads}

Look closely at a road on a printed map or a map app. It is not one line
but two: a darker, wider line underneath, called the \emph{casing}, and
the road's own colour on top. The casing gives the road an outline, which
makes it stand out from the land. And because every road's casing is drawn
before any road's colour, where two roads meet the colour of one paints
over the casing of the other, and the junction looks joined, not crossed.

\begin{center}
\begin{tikzpicture}[x=1cm,y=1cm]
  \begin{scope}
    \draw[line width=7pt, black!45] (0,0) -- (3,0);
    \draw[line width=5pt, white] (0,0) -- (3,0);
    \draw[line width=7pt, black!45] (1.5,-1) -- (1.5,1);
    \draw[line width=5pt, white] (1.5,-1) -- (1.5,1);
    \node[font=\small, text=cm] at (1.5,-1.4) {one road at a time};
  \end{scope}
  \begin{scope}[xshift=5cm]
    \draw[line width=7pt, black!45] (0,0) -- (3,0);
    \draw[line width=7pt, black!45] (1.5,-1) -- (1.5,1);
    \draw[line width=5pt, white] (0,0) -- (3,0);
    \draw[line width=5pt, white] (1.5,-1) -- (1.5,1);
    \node[font=\small, text=cm] at (1.5,-1.4) {all casings first};
  \end{scope}
\end{tikzpicture}
\end{center}

Tracks are drawn as a thin brown line with no casing, and footpaths as a
line of red dots, which is how most maps show them. One-way streets get
small arrows, but only when the map is zoomed in far enough for them to
fit.

There is one more thing to know about the canvas. Changing its pen, its
colour or width, takes far longer than drawing a short line: each change
goes through several calls into the Java libraries underneath. On a map of
a few hundred road segments, setting the pen once per road, twice over,
takes longer than all the drawing. So each pass draws all the roads of
one kind together, and sets the pen only once per kind.

\begin{continued}{map_view.nex}
  -- Roads are drawn in two passes, the way printed maps do it: first
  -- every road a little wider in a darker colour (the casing), then every
  -- road again in its own colour. Where roads meet, the second pass paints
  -- over the first, so junctions join cleanly.
  --
  -- Each pass draws the roads one kind at a time. Changing the pen's
  -- colour or width takes far longer than drawing a short line, so it
  -- pays to change it once per kind instead of once per road.
  draw_roads() do
    across [Road_Kind.Street, Road_Kind.Lane, Road_Kind.Main] as kind do
      set_pen(when kind = Road_Kind.Main then Palette.MAIN_CASING else Palette.CASING end,
        road_width(kind) + 2.0)
      across island.roads as r do
        if r.kind = kind then
          segments(island.points_of(r.node_ids))
        end
      end
    end
    across Road_Kind.values.reverse as kind do
      if kind = Road_Kind.Footpath then
        canvas.set_pen_color(Palette.FOOTPATH)
      elseif kind = Road_Kind.Track then
        set_pen(Palette.TRACK, 2.0)
      else
        set_pen(when kind = Road_Kind.Main then Palette.MAIN_ROAD else Palette.MINOR_ROAD end,
          road_width(kind))
      end
      across island.roads as r do
        if r.kind = kind then
          let pts: Array[Point] := island.points_of(r.node_ids)
          if kind = Road_Kind.Footpath then
            dots(pts, 1.3, 8.0)
          else
            segments(pts)
          end
        end
      end
    end
    if metres_per_pixel < 2.5 then
      canvas.set_pen_color(Palette.ROAD_LABEL)
      across island.roads as r do
        if r.is_one_way then
          draw_one_way_arrows(island.points_of(r.node_ids))
        end
      end
    end
  end

  -- The width in pixels of a road of a given kind.
  road_width(kind: Road_Kind): Real do
    if kind = Road_Kind.Main then
      result := 5.0
    elseif kind = Road_Kind.Lane then
      result := 4.0
    elseif kind = Road_Kind.Street then
      result := 3.5
    else
      result := 2.0
    end
  end

  -- Small arrowheads along a one-way road, pointing the way traffic goes.
  draw_one_way_arrows(pts: Array[Point]) do
    from
      let i: Integer := 1
    until
      i >= pts.length
    do
      let a: Point := pts.get(i - 1)
      let b: Point := pts.get(i)
      if not off_canvas(a, b) then
        arrowhead(a + (b - a) * 0.5, bearing(a, b), 5.0)
      end
      i := i + 1
    end
  end
\end{continued}

The second pass goes through the kinds in reverse, from footpaths to main
roads, so that more important roads are drawn on top of less important
ones.

Places get a dot: a white one with a smaller dark one on top, so that it
shows up on any background.

\begin{continued}{map_view.nex}
  -- A dot for every place: white, with a smaller dark dot on top.
  draw_places() do
    canvas.set_pen_color(Color.WHITE)
    across island.places as p do
      let at: Point := island.point(p.node_id)
      canvas.filled_circle(at.x, at.y, metres(5.0))
    end
    canvas.set_pen_color(Palette.LABEL)
    across island.places as p do
      let at: Point := island.point(p.node_id)
      canvas.filled_circle(at.x, at.y, metres(3.5))
    end
  end
\end{continued}

\section{The scale bar and the north arrow}

Two things every map should have: a \emph{scale bar}, to show how far
things are, and a \emph{north arrow}, to show which way is up. Our maps are
always drawn with north up, but saying so costs nothing.

A scale bar should be a round length: 500 m, not 487 m. We aim for about
150 pixels, and take the largest round length that fits.

\begin{continued}{map_view.nex}
  -- A bar of a round length (100 m, 200 m, 500 m, 1 km, ...) in the
  -- bottom left corner.
  draw_scale_bar() do
    let target: Real := metres(150.0)
    let length: Real := 10.0
    across [20.0, 50.0, 100.0, 200.0, 500.0, 1000.0, 2000.0, 5000.0] as candidate do
      if candidate <= target then
        length := candidate
      end
    end
    let left: Real := canvas.x_min + metres(20.0)
    let base: Real := canvas.y_min + metres(inset_bottom + 20.0)
    canvas.set_pen_color(Palette.HALO)
    canvas.filled_rectangle(left + length / 2.0, base + metres(9.0),
      length / 2.0 + metres(8.0), metres(16.0))
    canvas.set_pen_color(Palette.LABEL)
    canvas.set_pen_radius(1.0)
    canvas.line(left, base, left + length, base)
    canvas.line(left, base - metres(3.0), left, base + metres(3.0))
    canvas.line(left + length, base - metres(3.0), left + length, base + metres(3.0))
    canvas.set_font(Font.SANS.with_size(11))
    canvas.text(left + length / 2.0, base + metres(12.0), distance_text(length))
  end

  -- An arrow pointing north, in the top right corner.
  draw_north_arrow() do
    let x: Real := canvas.x_max - metres(30.0)
    let y: Real := canvas.y_max - metres(inset_top + 45.0)
    canvas.set_pen_color(Palette.LABEL)
    canvas.filled_polygon([x, x - metres(8.0), x + metres(8.0)],
      [y + metres(14.0), y - metres(10.0), y - metres(10.0)])
    canvas.set_font(Font.SANS.bold().with_size(13))
    canvas.text(x, y + metres(26.0), "N")
  end
\end{continued}

The canvas's \code{x_min}, \code{x_max}, \code{y_min} and \code{y_max} are
the edges of what it shows, so these are placed relative to the corners,
whatever the zoom. \code{filled_rectangle} takes a centre and half the
width and height, and the canvas's \code{filled_polygon} takes the corners'
$x$ and $y$ coordinates as two separate arrays.

\section{Drawing tools}

Last, the drawing tools used above, which the live map will use too, for
drawing routes and markers on top of the map.

A line through a list of points is drawn one segment at a time. Before
drawing a segment, we check whether it is anywhere near the canvas. When
the map is zoomed in, most of the island is off the canvas, and drawing it
would be wasted work. A segment is certainly off the canvas if both its
ends are beyond the same edge, which takes a few comparisons to check.
(A segment can be off the canvas without that, cutting across a corner,
but drawing a few of those is harmless.)

\begin{continued}{map_view.nex}
  -- Drawing tools, used by the map and by overlays such as a route.

  -- A line through `pts`, `width` pixels wide.
  line(pts: Array[Point], color: Color, width: Real) do
    set_pen(color, width)
    segments(pts)
  end

  set_pen(color: Color, width: Real) do
    canvas.set_pen_color(color)
    canvas.set_pen_radius(width / 2.0)
  end

  -- The line through `pts`, with the pen as it is.
  segments(pts: Array[Point]) do
    from
      let i: Integer := 1
    until
      i >= pts.length
    do
      let a: Point := pts.get(i - 1)
      let b: Point := pts.get(i)
      if not off_canvas(a, b) then
        canvas.line(a.x, a.y, b.x, b.y)
      end
      i := i + 1
    end
  end

  -- Whether the segment a-b is certainly off the canvas, because both
  -- its ends are beyond the same edge. Drawing it would be wasted work.
  off_canvas(a, b: Point): Boolean do
    let m: Real := metres(12.0)
    result := (a.x < canvas.x_min - m and b.x < canvas.x_min - m)
      or (a.x > canvas.x_max + m and b.x > canvas.x_max + m)
      or (a.y < canvas.y_min - m and b.y < canvas.y_min - m)
      or (a.y > canvas.y_max + m and b.y > canvas.y_max + m)
  end
\end{continued}

The canvas's pen \emph{radius} is half the line's width, hence
\code{width / 2.0} in \code{set_pen}.

A dotted line puts a dot every \code{gap} pixels along the line. The dots
must be evenly spaced along the whole line, not restarting at each segment,
so we keep track of how far along the whole line the current segment
starts, and where along the line the next dot is due.

\begin{continued}{map_view.nex}
  -- Dots along the line through `pts`, in the pen's colour.
  dots(pts: Array[Point], radius, gap: Real) do
    let step: Real := metres(gap)
    -- How far along the whole line the current segment starts, and where
    -- the next dot goes.
    let start: Real := 0.0
    let next_dot: Real := 0.0
    from
      let i: Integer := 1
    until
      i >= pts.length
    do
      let a: Point := pts.get(i - 1)
      let b: Point := pts.get(i)
      let finish: Real := start + a.distance_to(b)
      if off_canvas(a, b) then
        -- Skip the dots this segment would have had.
        if next_dot <= finish then
          next_dot := next_dot + ((finish - next_dot) / step).floor.to_real * step + step
        end
      else
        from
        until
          next_dot > finish
        do
          let p: Point := a + (b - a) * ((next_dot - start) / (finish - start))
          canvas.filled_circle(p.x, p.y, metres(radius))
          next_dot := next_dot + step
        end
      end
      start := finish
      i := i + 1
    end
  end
\end{continued}

A filled polygon needs its corners as two arrays, and an arrowhead is a
small triangle: its tip, and two corners a little way back from the tip,
one on each side. \code{direction(heading)} is an arrow of length 1
pointing along the heading, and \code{direction(heading + 90.0)} points
square to its right.

\begin{continued}{map_view.nex}
  fill_polygon(pts: Array[Point]) do
    let xs: Array[Real] := []
    let ys: Array[Real] := []
    across pts as p do
      xs.add(p.x)
      ys.add(p.y)
    end
    canvas.filled_polygon(xs, ys)
  end

  -- A filled triangle `size` pixels long at `tip`, pointing along
  -- `heading` (a compass bearing).
  arrowhead(tip: Point, heading: Real, size: Real) do
    let forward: Point := direction(heading) * metres(size)
    let side: Point := direction(heading + 90.0) * metres(size * 0.6)
    let back: Point := tip - forward
    canvas.filled_polygon([tip.x, back.x + side.x, back.x - side.x],
      [tip.y, back.y + side.y, back.y - side.y])
  end

  show() do
    canvas.show()
  end

  save(file: String) do
    canvas.save(user_path(file))
  end

private feature
  area_color(kind: String): Color do
    if kind = "water" then
      result := Palette.WATER
    elseif kind = "wood" then
      result := Palette.WOOD
    elseif kind = "sand" then
      result := Palette.SAND
    elseif kind = "town" then
      result := Palette.TOWN
    else
      result := Palette.LAND
    end
  end
end
\end{continued}

\section{The first picture}

\begin{tryit}{try_draw.nex}
intern Map_View

let island: Island_Map := create Island_Map.load(user_path("kestrel_island.map"))
let view: Map_View := create Map_View.make(island, 1000, 700)
view.draw_map()
view.save("first_map.png")
say("saved first_map.png, at " + fixed(view.metres_per_pixel, 2) + " metres per pixel")
\end{tryit}
\tryoutput{try_draw}
\mappicture{first_map}

That is Kestrel Island. The orange ring is the Coast Road, the white lines
are lanes and the town's streets, the brown lines are tracks through the
woods, and the red dots are footpaths. The blue line from the lake to the
sea is Kestrel Brook. Port Kestrel is the grid of streets on the south
coast.

\section{Why we projected}

What if we had drawn the latitudes and longitudes directly, a degree east
the same size as a degree north? This script does exactly that for the
island and its lake, scaling degrees to metres with the length of a degree
of latitude, both ways:

\begin{tryit}{try_unprojected.nex}
intern Map_View

-- A node's position, pretending a degree is the same length both ways.
function naive(ll: Lat_Lon, middle: Lat_Lon, scale: Real): Point do
  result := create Point.make((ll.longitude - middle.longitude) * scale,
    (ll.latitude - middle.latitude) * scale)
end

let island: Island_Map := create Island_Map.load(user_path("kestrel_island.map"))
let view: Map_View := create Map_View.make(island, 1000, 700)
let k: Real := island.projection.metres_per_degree_lat
let centre: Lat_Lon := island.projection.centre
view.look_at(create Point.make(0.0, 0.0), view.metres_per_pixel * 1.45)
view.canvas.clear(Palette.SEA)
across island.areas as a do
  if a.kind = "land" or a.kind = "water" then
    let pts: Array[Point] := []
    across a.node_ids as id do
      pts.add(naive(island.positions.get(id), centre, k))
    end
    view.canvas.set_pen_color(when a.kind = "land" then Palette.LAND else Palette.WATER end)
    view.fill_polygon(pts)
  end
end
view.save("unprojected.png")
say("saved unprojected.png")
\end{tryit}
\tryoutput{try_unprojected}
\mappicture[0.8\linewidth]{unprojected}

The island is stretched sideways by $1/\cos 47^\circ$, about one and a
half times, and the round lake has become a long oval. Every distance and
every angle measured on that picture would be wrong.
